\CourseSection{Introduction}\label{sec:introduction}
\Guidance{The body of the completed proposal, \emph{i.e.} Sections 1--6 below, has a maximum length of eight pages; the title page, executive summary, contents, references, and appendices are not included in this limit. \\
In the intro below, establish the wider problem you are trying to address, project's scope in the context of this problem, the general design, and limitations. \\
Change the titles of these subsections as you see fit. We have just provided some example placeholders.}

\CourseSubsection{Challenges in Golf Course Navigation and Player Fatigue}\label{sec:problem-and-purpose}

\Guidance{Describe the problem, the intended end user or market, evidence of the need, and the value of addressing it. Support the context with several citations from peer-reviewed literature. Feel free to change the section titles to fit your project. This section is the 'Problem and Purpose' section.}
The problem this project addresses is the physicality and brain of walking on a golf course. The intended end user is more toward a senior golfer, who is often burdened by the physicality of walking on a golf course and the cognitive load of knowing which club to use and not leaving clubs behind. The value of addressing this problem is that it will allow senior golfers to enjoy the game more, and potentially improve their performance by reducing fatigue and cognitive load.

\CourseSubsection{Smart Autonomous Caddy Design for Physical and Cognitive Offloading}\label{sec:project-objectives}
\Guidance{State the specific objectives and measurable outcomes that address the problem. Characterize the electrical and / or computer engineering involved. Support your claims by describing similar approaches in the peer-reviewed literature. Include a general diagram showing the system context. Compare your approach to those in the peer-reviewed literature. This section is the 'Project Objective' section.}
The primary objective of this project is to design and build a functional prototype of an autonomous, motorized golf caddy system that seamlessly tracks and follows a golfer, assists with club selection, and prevents misplaced equipment. By integrating computer vision, embedded sensor fusion, and power-efficient motor control, the system eliminates the physical strain of manually pulling or carrying bags over varied terrain.
To address the specific needs of senior golfers, the engineering scope encompasses three key engineering domains:
\begin{enumerate}
\item \textbf{Autonomous Navigation \& Following:} A computer-vision and ultra-wideband (UWB) tracking subsystem running on an embedded microcontroller to maintain a safe follow distance, avoid field obstacles, and eliminate manual steering effort.
\item \textbf{Cognitive Assisting \& Equipment Management:} An RFID/vision-based club tracking system combined with a basic distance-recommendation algorithm to track which clubs are removed and ensure none are left behind on the green.
\item \textbf{Mechatronic Power System:} A high-torque DC motor drive architecture optimized for regenerative braking and hill traversal without tipping or uncontrolled acceleration on slopes.
\end{enumerate}

\begin{figure}[htbp]
\centering
% Placeholder for your system diagram
\rule{0.8\textwidth}{0.3\textwidth}
\caption{System context diagram illustrating the interaction between the user location beacon, embedded vision processing, motor actuation, and club tracking subsystems.}
\label{fig:system_context}
\end{figure}

Existing commercial solutions rely heavily on expensive, fully autonomous robotic chassis or basic remote-controlled cart motors that lack intelligence and automatic club-tracking. Our approach integrates localized sensor fusion directly into a retrofitted lightweight chassis, achieving lower deployment cost while directly targeting cognitive relief alongside mobility assistance.

\CourseSubsection{Scope and Deliverables: Autonomous Tracking, Smart Club Auditing, and Safety Controls}\label{sec:project-scope}
\Guidance{This is a refinement of Sec. \ref{sec:project-objectives}. Define the end product or products, the operating context, system inputs and outputs, and the major components. Explain each component's purpose and how the components fit together. State the project boundaries and assumptions. Address several limitations of your approach. Provide citations in the peer-reviewed literature and / or textbooks that apply similar strategies. This section is the 'Scope and deliverables' section.}
The primary deliverable for this project is a fully integrated, low-power smart golf caddy prototype capable of operating across standard 18-hole terrain under typical outdoor environmental conditions.
\subsubsection*{System Boundary and Operating Context}
The system operates within the physical boundary of standard outdoor golf course terrain (grass, fairways, mild incline pathways up to $15^\circ$). Inputs to the system include real-time distance/angle data from user tracking beacons, wheel encoder feedback, and manual override controls via a handlebar/remote toggle. Outputs consist of differential motor drive power signals, audio/visual alerts for unreturned clubs, and course/distance readouts displayed on an onboard user interface.

\subsubsection*{Major System Components}
\begin{itemize}
\item \textbf{Main Controller Unit (MCU/SBC):} Coordinates sensor telemetry, executes path-following logic, and monitors battery state of charge.
\item \textbf{Sensor Array (UWB + Ultrasonic/LiDAR):} Measures precise range and bearing to the user while detecting ground obstacles and steep drop-offs.
\item \textbf{Club Inventory Subsystem:} Uses short-range RFID tags on club grips to audit active equipment status in real time.
\item \textbf{Motor Controller and Drive System:} Dual-channel MOSFET motor driver with closed-loop PID velocity control to ensure smooth acceleration on variable grass textures.
\end{itemize}

\subsubsection*{Limitations and Assumptions}
\begin{itemize}
\item \textbf{Terrain Constraints:} The prototype assumes operation on standard course turf and pathways; extreme mud, deep sand traps, and water hazards are outside the current operational scope.
\item \textbf{Environmental Conditions:} Optical/camera tracking accuracy may degrade in heavy rain or dense fog, requiring the system to fall back on low-frequency RF rangefinding.
\item \textbf{Payload Limit:} The chassis is rated to transport standard bag weights up to $15\text{ kg}$, excluding passenger ride-along capabilities.
\end{itemize}

\clearpage
